\section{Matrizes}
As matrizes representam indivíduos ou sujeitos nas linhas e variáveis nas colunas. Ou seja, coletarmos 4 informações de 10 indivíduos, teremos uma matriz com 10 linhas e 4 colunas. As as 4 informações de cada indivíduo forem representas em um vetor coluna, para construir a matriz de dados, podemos empilhar em cada linha as transpostas desses  vetores colunas.

\begin{equation}
\begin{aligned}
X = \begin{bmatrix}
    \mbox{---}  \mathbf{x}^T_1 \mbox{---}  \\
     \mbox{---}  \mathbf{x}^T_2 \mbox{---}  \\
    \vdots \\
     \mbox{---}  \mathbf{x}^T_n \mbox{---}  \\
\end{bmatrix} 
\end{aligned}
\end{equation}